% ===== Appendix: corrections to the book, what we added, sources (after Questions; not presented) =====
\appendix
\partaccent{lavender}\gdef\idiottopicemoji{books}\gdef\idiottopicname{Հավելված}

\begin{frame}{Հավելված Ա. Ուղղումներ գրքի նկատմամբ (1/2)}
  \scriptsize
  \setlength{\tabcolsep}{4pt}
  \renewcommand{\arraystretch}{1.15}
  \begin{tabular}{@{}L{2.4cm}L{3.4cm}L{4.5cm}L{3.2cm}@{}}
    \toprule
    \textbf{Թեմա} & \textbf{Գրքում (2016)} & \textbf{Իրականում} & \textbf{Աղբյուր}\\
    \midrule
    Գեյջի վնասվածքը & Ձողը ոչնչացրեց ձախ ճակատային բլթի «գրեթե ամբողջը» & Կեղևի մոտ 4\%-ը, բայց սպիտակ նյութի 10\%-ից ավելին (կապերը) & Van Horn et al., PLoS ONE, 2012\\
    Գծերի մասին մեջբերումը & «Ջ. Պ. Գիլարդ, 1959» & Ջ. Պ. Գիլֆորդ (Guilford), «Personality», 1959 & Guilford, 1959\\
    Երկվորյակների թիվը & «Գրեթե 300 զույգ» & 250 զույգ (123 միաձվային + 127 երկձվային) & Jang et al., J. Pers., 1996\\
    Ա և Բ տիպեր & «Ֆրիդման ու Ռոզենհան» & Ֆրիդման ու Ռոզենման (Rosenman). գրքի ծանոթագրությունում ճիշտ է & Friedman \& Rosenman, 1974\\
    MBTI & Myers-Briggs Type «Inventory» & Myers-Briggs Type Indicator & Myers \& Briggs\\
    Բառային վարկած & «Բառերը կարող են լինել այնքան, որքան գծեր կան» & Կարևոր տարբերությունները լեզվում բառեր են դառնում. որքան կարևոր, այնքան շատ բառ & Goldberg, Am. Psychol., 1993\\
    \bottomrule
  \end{tabular}
\end{frame}

\begin{frame}{Հավելված Ա. Ուղղումներ գրքի նկատմամբ (2/2)}
  \scriptsize
  \setlength{\tabcolsep}{4pt}
  \renewcommand{\arraystretch}{1.15}
  \begin{tabular}{@{}L{2.4cm}L{3.4cm}L{4.5cm}L{3.2cm}@{}}
    \toprule
    \textbf{Թեմա} & \textbf{Գրքում (2016)} & \textbf{Իրականում} & \textbf{Աղբյուր}\\
    \midrule
    Վտանգի համակարգը & Նշաձև մարմին, հիպոկամպ, գորշ նյութ՝ «միջին ուղեղի մասեր» & Միջին ուղեղում միայն գորշ նյութն է, մյուսները քունքային բլթում են & Դասագրքային\\
    Զայրույթը «թափել» & Մի զսպիր, թող դուրս գա & «Թափելը» (venting) չի օգնում, հանդարտեցումն է օգնում & Kjærvik \& Bushman, 2024\\
    Մասլոուի բուրգ & Աստիճաններ՝ խիստ հերթականությամբ & Մասլոուն բուրգ չի նկարել. հերթականությունը «այդքան կոշտ չէ» & Maslow, 1943; Bridgman et al., 2019\\
    Ինքնորոշման տեսություն & 1988, «կապվածություն»՝ ճանաչում & 1985, կապվածություն՝ կապ ուրիշների հետ & Ryan \& Deci, 2000\\
    Զեյգարնիկի էֆեկտ & Անավարտը ավելի լավ է հիշվում & Չի հաստատվում. վերադառնալու հակումը կա & Ghibellini \& Meier, 2025\\
    Գորտի մեջբերումը & «Է. Լ. Ուայթ» & Է. Բ. Ուայթ ու Քեթրին Ուայթ, 1941 & Quote Investigator, 2014\\
    \bottomrule
  \end{tabular}
\end{frame}

\begin{frame}{Հավելված Բ. Ինչ ենք ավելացրել գրքից դուրս}
  \begin{columns}[T]
    \begin{column}{0.48\textwidth}
      \begin{card}[colback=mint!30!cardwhite, fontupper=\scriptsize]
        {\blackfont Թվեր ու հետազոտություններ}\par\smallskip
        \begin{itemize}\setlength{\itemsep}{0.3mm}
          \item Գեյջի վնասի չափը՝ Van Horn et al., 2012
          \item Ֆրենոլոգ Սայզերի պնդումը՝ Macmillan, 2000
          \item Ժառանգականության թվերը՝ Jang et al., 1996
          \item MBTI-ի կրկնակի թեստը՝ Pittenger, 1993, 2005
          \item Ֆլոմաստերների փորձի թվերը՝ Lepper et al., 1973
          \item Venting՝ Kjærvik \& Bushman, 2024
          \item Մասլոուի բուրգը՝ Bridgman et al., 2019
          \item Զեյգարնիկ՝ Ghibellini \& Meier, 2025
        \end{itemize}
      \end{card}
    \end{column}
    \begin{column}{0.48\textwidth}
      \begin{card}[colback=peach!35!cardwhite, fontupper=\scriptsize]
        {\blackfont Մերն է, ոչ թե գրքից}\par\smallskip
        \begin{itemize}\setlength{\itemsep}{0.3mm}
          \item Օնլայն թեստերի օրինակները
          \item Ծխնելույզի կատակը որպես «անավարտ գործ»
          \item Վարագույրի կատակի պուանտը ու ձիու կատակի բացատրությունը
          \item Զայրույթ-թշնամանք-ագրեսիա վիկտորինան
          \item Երկվորյակների մեթոդի բացատրությունը
          \item Հայերեն բառերի չաթը
          \item Մեմերը. բովանդակությունը գրքից է
        \end{itemize}
      \end{card}
    \end{column}
  \end{columns}
\end{frame}

\begin{frame}{Հավելված Գ. Աղբյուրներ (1/2)}
  \scriptsize
  Գրքի ծանոթագրություններից, որոնց վրա հիմնված են սլայդները\par\smallskip
  \begin{columns}[T]
    \begin{column}{0.48\textwidth}
      \begin{itemize}\setlength{\itemsep}{0.4mm}
        \item Phares E. J., Chaplin W. F. Introduction to Personality. 1997
        \item Froman L. A. Personality and political socialization. J. Politics, 1961
        \item Eysenck H., Levey A. Conditioning, introversion-extraversion and the strength of the nervous system. 1972
        \item Taki Y. et al. Personality traits and regional gray matter volume changes. Hum. Brain Mapp., 2013
        \item Jang K. L., Livesley W. J., Vernon P. A. Heritability of the Big Five personality dimensions. J. Pers., 1996
        \item Friedman M., Rosenman R. H. Type A Behavior and Your Heart. 1974
        \item Caprara G. V., Cervone D. Personality: Determinants, Dynamics, and Potentials. 2000
      \end{itemize}
    \end{column}
    \begin{column}{0.48\textwidth}
      \begin{itemize}\setlength{\itemsep}{0.4mm}
        \item Murray J. B. Review of research on the Myers-Briggs Type Indicator. Percept. Mot. Skills, 1990
        \item Sell A. N. The recalibrational theory and violent anger. Aggress. Violent Behav., 2011
        \item Carver C. S., Harmon-Jones E. Anger is an approach-related affect. Psychol. Bull., 2009
        \item Kazén M. et al. Inverse relation between cortisol and anger. Biol. Psychol., 2012
        \item Rutherford H. J., Lindell A. K. Thriving and surviving: approach and avoidance motivation and lateralization. Emot. Rev., 2011
        \item Antos D. et al. Emotion expression and trustworthiness in negotiation. AAAI, 2011
      \end{itemize}
    \end{column}
  \end{columns}
\end{frame}

\begin{frame}{Հավելված Գ. Աղբյուրներ (2/2)}
  \scriptsize
  Գրքի ծանոթագրություններից, որոնց վրա հիմնված են սլայդները\par\smallskip
  \begin{columns}[T]
    \begin{column}{0.48\textwidth}
      \begin{itemize}\setlength{\itemsep}{0.4mm}
        \item Freud S. Beyond the Pleasure Principle. 1920
        \item McLeod S. Maslow's hierarchy of needs. Simply Psychology, 2014
        \item Ryan R. M., Deci E. L. Self-determination theory. Am. Psychol., 2000
        \item Lepper M. R., Greene D., Nisbett R. E. Undermining children's intrinsic interest with extrinsic reward. JPSP, 1973
        \item Higgins E. T. Self-discrepancy: a theory relating self and affect. Psychol. Rev., 1987
        \item Reeve J., Cole S. G., Olson B. C. The Zeigarnik effect and intrinsic motivation. Motiv. Emot., 1986
        \item Shuster S. Sex, aggression, and humour: responses to unicycling. BMJ, 2007
      \end{itemize}
    \end{column}
    \begin{column}{0.48\textwidth}
      \begin{itemize}\setlength{\itemsep}{0.4mm}
        \item Bell N. D. Responses to failed humor. J. Pragmatics, 2009
        \item Shurcliff A. Judged humor, arousal, and the relief theory. JPSP, 1968
        \item Hayworth D. The social origin and function of laughter. Psychol. Rev., 1928
        \item Provine R. R., Emmorey K. Laughter among deaf signers. J. Deaf Stud. Deaf Educ., 2006
        \item Provine R. R. Contagious laughter. Bull. Psychon. Soc., 1992
        \item McGettigan C. et al. Individual differences in laughter perception. Cereb. Cortex, 2015
      \end{itemize}
    \end{column}
  \end{columns}
\end{frame}
